\subsection{Dirac-Delta Distribution $\delta(t)$}
\begin{emphbox}
    \begin{equation*}
        \int_{-\infty}^\infty f(t)\delta(t - a)\diff t = f(a) \quad \forall f:\mathbb R \rightarrow \mathbb C \text{, die bei } a \text{ stetig}
    \end{equation*}
\end{emphbox}

\begin{sectionbox}
    \subsubsection{$\delta$-Folge}
    Folge an Funktionen, für die gilt
    \begin{itemize}
        \item[] \begin{equation*}
                  \int_{-\infty}^{\infty}\delta_n(t)\diff t = 1 \; \forall n \in \mathbb Z
              \end{equation*}
        \item[] \begin{equation*}
                  \lim_{n\rightarrow\infty}\int_{-\infty}^{\infty}f(t)\delta_n(t)\diff t = f(0)
              \end{equation*} für jede in $0$ stetige Funktion $f$
    \end{itemize}
\end{sectionbox}

\subsection{Heaviside Theta/Stufenfunktion}
\begin{sectionbox}
    \begin{equation*}
        \theta(t) = \begin{cases}
            0           & t < 0 \\
            \frac{1}{2} & t = 0 \\
            1           & t > 0
        \end{cases} = \lim_{n\rightarrow\infty}\theta_n(t)
    \end{equation*}
    mit z.B. \begin{equation*}
        \theta_n(t)=\frac{1}{2}(\tanh(tn) + 1) = \frac{1}{2}\text{sign}(t)e^{-|t|/n} + \frac{1}{2}
    \end{equation*}
\end{sectionbox}

\begin{sectionbox}
    \subsection{Bekannte Fouriertransformationen}
    \begin{itemize}
        \item $\hat \delta(\omega) = 1(\omega) \qquad \hat 1 = 2\pi\delta$
    \end{itemize}
\end{sectionbox}

\subsection{Verallgemeinerte Ableitung}
\begin{emphbox}
    Sei $f:\mathbb R \rightarrow \mathbb C$, wobei 
    \begin{itemize}
        \item $f'(t)$ für alle $t\in\mathbb R\setminus \{t_1 < ... < t_N\}$ existiert ($f'(t_j)=0$)
        \item Die einseitigen Grenzwerte für alle Unstetigkeitsstellen existieren
    \end{itemize}
    Dann gilt
    \begin{equation*}
        Df(t) = f'(t) + \sum_{j=1}^{N}(f(t_j^+) - f(t_j^-))\delta(t-t_j)
    \end{equation*}
    mit Spektralfunktion $\widehat{Df}(\omega) = i\omega \hat f(\omega)$
\end{emphbox}

\subsection{Abtasttheorem}
Sei $f\in L^1(\mathbb R)$ stückweise stetig diffb. und bandbegrenzt ($\hat f(\omega) = 0 \; \forall \omega >\Omega$). 
Sei $\omega_s=\frac{2\pi}{\Delta t}$ eine Abtastfrequenz mit $\omega_s > 2\Omega$. Dann gilt an jeder Stetigkeitsstelle $t$ von $f$, dass $f$ mittels
\begin{equation*}
    f(t)=\sum_{k\in\mathbb Z}f(k\Delta t)\sinc\left[\frac{\pi}{\Delta t}(t-k\Delta t)\right]
\end{equation*} rekonstruiert werden kann.
Mit $\sinc x=\sin(x)/x$